\subsection{Trigger Nodes Mechanism}

\subsubsection{Overview}
In a workflow automation platform, every workflow must have an entry point — something that decides when execution begins. This role is fulfilled by trigger nodes. A trigger node is a specialized node that listens for a specific event or condition and, when that condition is met, fires the rest of the workflow graph for execution.
Unlike regular action nodes that simply process inputs and pass outputs downstream, trigger nodes are always-on components that run persistently as long as a workflow is deployed. Each trigger type reacts to a different source of events. The platform supports the following trigger types:
\begin{itemize}
  \item \textbf{Cron} --- fires repeatedly on a fixed time interval, starting from a configurable date.
  \item \textbf{Poll} --- periodically checks an external source for new data.
  \item \textbf{Webhook} --- waits for an incoming HTTP request before firing.
  \item \textbf{Manual} --- fires immediately on demand when a user triggers execution directly.
\end{itemize}

The trigger mechanism is built around three key concerns: starting triggers when a workflow is deployed, reliably passing execution down the workflow graph when a trigger fires, and cleanly stopping all triggers when a workflow is stopped or the server shuts down.

\subsubsection{Core Abstractions}
\paragraph{The EventSource Interface}
All trigger types implement the \texttt{EventSource} interface, which defines a minimal two-method contract:
\begin{lstlisting}
// eventSource.ts
import { Job } from 'bullmq';
 
export interface EventSource {
  start(body: Record<string, unknown>): Promise<Job>;
}
\end{lstlisting}

This interface was intentionally kept minimal. The \texttt{start()} method receives the node configuration and returns a BullMQ Job handle, while \texttt{stop()} is responsible for cleanly cancelling whatever mechanism the trigger started. Any new trigger type only needs to implement these two methods, making the system open to extension without modifying existing code.
 
Enforcing the same \texttt{start} contract across all trigger types means the workflow runner can manage cron and poll triggers through a single loop, without any conditional branching based on type. Adding a new trigger type in the future requires no changes to the runner --- only a new class implementing \texttt{EventSource}.

\paragraph{The Node Base Class}
Trigger nodes are still Node subclasses, following the same abstract base class as all other node types in the system. What distinguishes a trigger node is the \texttt{triggerNode: true} flag in its static description object, along with a type of 'cron', 'poll', etc.

\begin{lstlisting}
// nodeBase.ts (excerpt)
export interface INodeDescription {
  name: string;
  displayName: string;
  type: 'cron' | 'curl' | 'webhook' | 'generic' | 'poll';
  category: 'ai' | 'action-in-app' | 'data-transformation' | 'core' | 'flow-control';
  inputs: string[];
  outputs: string[];
  credentials?: string[];
  parameters?: any[];
  triggerNode?: boolean;   // marks this node as a trigger
}
 
export abstract class Node {
  static description: INodeDescription;
  abstract execute(inputs: any): Promise<any>;
  poll?(): Promise<any>;
  webhook?(request: any): Promise<any>;
}
\end{lstlisting}

The \texttt{execute()} method is abstract and must be implemented by every node, including trigger nodes. For trigger nodes, this method represents what happens when the trigger actually fires — for a cron node it simply returns a confirmation string, since its real work is handled by the scheduler.
